\documentclass[conference]{IEEEtran}
\IEEEoverridecommandlockouts

% Essential Packages
\usepackage{cite}
\usepackage{amsmath,amssymb,amsfonts}
\usepackage{algorithmic}
\usepackage{graphicx}
\usepackage{textcomp}
\usepackage{xcolor}
\usepackage{booktabs}
\usepackage{multirow}
\usepackage{url}
\usepackage{microtype}
\usepackage{subcaption}

\def\BibTeX{{\rm B\kern-.05em{\sc i\kern-.025em b}\kern-.08em
    T\kern-.1667em\lower.7ex\hbox{E}\kern-.125emX}}


\usepackage{enumitem}
\setlist{nosep,leftmargin=*}
\setlength{\textfloatsep}{8pt plus 2pt minus 2pt}
\setlength{\floatsep}{6pt plus 2pt minus 2pt}
\setlength{\intextsep}{6pt plus 2pt minus 2pt}
\setlength{\abovedisplayskip}{3pt}
\setlength{\belowdisplayskip}{3pt}

\begin{document}

\title{Curvature-Aware Directional Swift Attention Network (CAD-SPAN) for Efficient Single Image Super-Resolution}

\author{
\IEEEauthorblockN{Aniket Gupta}
\IEEEauthorblockA{\textit{Department of Information Technology} \\
\textit{Indian Institute of Information Technology, Allahabad}\\
Prayagraj, India \\
iit2023018@iiita.ac.in}
\and
\IEEEauthorblockN{Ritesh Narayan Das}
\IEEEauthorblockA{\textit{Department of Information Technology} \\
\textit{Indian Institute of Information Technology, Allahabad}\\
Prayagraj, India \\
iit2023060@iiita.ac.in}
\and
\IEEEauthorblockN{Siddhartha Srivastava}
\IEEEauthorblockA{\textit{Department of Information Technology} \\
\textit{Indian Institute of Information Technology, Allahabad}\\
Prayagraj, India \\
iit2023504@iiita.ac.in}
}

\maketitle

\begin{abstract}
Single Image Super-Resolution (SISR) is an ill-posed low-level computer vision problem requiring the recovery of high-frequency spatial details from degraded, low-resolution inputs while conforming to strict edge-deployment constraints. While contemporary deep attention networks and Vision Transformers achieve competitive reconstruction fidelity, their heavy parameter footprints, massive memory access costs (MACs), and prolonged inference latencies severely inhibit edge and mobile deployment. The recently proposed Swift Parameter-free Attention Network (SPAN) established a new efficiency benchmark by employing an origin-symmetric activation function to derive attention maps directly from intermediate convolutional features without learnable parameters, winning the NTIRE 2024 Efficient Super-Resolution Challenge.

In this work, we first replicate the official SPAN baseline, validating reported reconstruction fidelity on standard benchmark datasets (achieving 32.19 dB PSNR on Set5, 28.63 dB on Set14, and 27.36 dB on BSD100 for scale factor $\times 4$). We then uncover three fundamental structural and theoretical limitations: (1) the origin-symmetric activation function is empirically formulated without a rigorous analytical foundation; (2) its point-wise formulation is strictly isotropic, treating all directional gradients identically; and (3) it entirely lacks differential curvature sensitivity along complex contours and junctions.

To overcome these deficiencies, we propose \textbf{CAD-SPAN} (Curvature-Aware Directional Swift Parameter-free Attention Network). First, we prove mathematically that SPAN's origin-symmetric activation is strictly equivalent to a scaled hyperbolic tangent gate ($\sigma(x) - 0.5 \equiv \frac{1}{2}\tanh(\frac{x}{2})$), and generalize it through temperature-calibrated steepness ($\tau = 0.502$). Second, grounded in differential geometry, we introduce \textbf{Curvature-Aware Directional Attention (CAD-SA)}, which decomposes feature representations into first-order directional gradient flux ($\|\nabla H\|_2$) via orthogonal Sobel operators and second-order isotropic curvature ($\Delta H$) via discrete Laplacian operators. Comprehensive empirical evaluations across Set5, Set14, and BSD100 confirm that CAD-SPAN consistently outperforms the NTIRE 2024 challenge winner in both PSNR and SSIM (up to $+0.078$ dB PSNR gain overall and up to $+0.160$ dB on high-frequency textured scenes) while strictly guaranteeing zero-degradation identity fallback.
\end{abstract}

\begin{IEEEkeywords}
Image Super-Resolution, Lightweight Attention Modules, Parameter-free Attention, Differential Geometry, Origin-Symmetric Activation, Curvature Prior, Re-parameterization.
\end{IEEEkeywords}

\section{Introduction and Related Work}

\subsection{Single Image Super-Resolution Fundamentals}
Single Image Super-Resolution (SISR) seeks to estimate a latent high-resolution (HR) image $I_{HR} \in \mathbb{R}^{H \times W \times C}$ from its low-resolution (LR) counterpart $I_{LR} \in \mathbb{R}^{\frac{H}{s} \times \frac{W}{s} \times C}$, where $s$ denotes the upscale factor. Because the downsampling operator is non-invertible, multiple feasible high-resolution images correspond to a single degraded input.

\subsection{Evolution of Attention in Super-Resolution}
Since the advent of deep learning in SISR initiated by SRCNN \cite{dong2015image}, restoration fidelity has been driven forward by deeper residual networks (e.g., EDSR \cite{lim2017enhanced}, RCAN \cite{zhang2018image}) and recent Vision Transformers (e.g., SwinIR \cite{liang2021swinir}, HAT \cite{chen2023activating}). A primary catalyst behind these advancements is the attention mechanism. Attention modules allow the neural network to dynamically allocate computational capacity toward high-frequency structural transitions—such as sharp edges, repetitive textures, and fine boundaries—where blurring and ringing artifacts predominantly emerge.

However, standard spatial and channel attention mechanisms (e.g., Squeeze-and-Excitation \cite{hu2018squeeze}, CBAM \cite{woo2018cbam}, Non-Local blocks \cite{wang2018non}) come with non-negligible computational overheads:
\begin{enumerate}
    \item \textbf{Substantial Parameter Overhead}: Learnable reduction/expansion projection convolutions inflate parameter footprints.
    \item \textbf{Memory Access Costs (MAC)}: Multi-branch tensor operations, global pooling reductions, and matrix cross-multiplications break memory locality and incur heavy GPU/NPU memory access latency.
    \item \textbf{Inference Latency Bottlenecks}: Complex routing branches impede layer fusion and hardware compiler optimizations, preventing real-time throughput on mobile and embedded devices.
\end{enumerate}

\begin{figure*}[t]
\centering
\includegraphics[width=0.96\textwidth,height=2.3in,keepaspectratio]{methodology_architecture.png}
\caption{\textbf{Comprehensive Methodology Architecture of Proposed CAD-SPAN}. (a) Macro-Architecture: Low-resolution input $I_{LR}$ is projected by a re-parameterizable $\text{Conv3XC}$ layer ($F_0$), processed through a cascade of 6 Curvature-Aware Directional Swift Parameter-free Attention Blocks (CAD-SPAB), aggregated via hierarchical multi-stage concatenation, and upscaled by sub-pixel convolution (PixelShuffle). (b) Micro-Architecture of CAD-SPAB: High-level feature $H$ is extracted via three consecutive $\text{Conv3XC}$ convolutions with SiLU activations. $H$ is simultaneously evaluated through: (1) Temperature-calibrated origin-symmetric gating $A_{\text{sym}} = \frac{1}{2}\tanh(\tau \cdot H)$; (2) First-order directional gradient flux $\|\nabla H\|_2$ via orthogonal Sobel kernels ($K_x, K_y$); and (3) Second-order isotropic curvature $\Delta H$ via discrete Laplacian kernel $K_L$. The resulting curvature-modulated gate $\widetilde{A}$ modulates the residual feature $(H + X_{\text{in}})$.}
\label{fig:methodology_arch}
\end{figure*}

\subsection{Efficient Super-Resolution Topologies}
Efficient SISR focuses on maximizing reconstruction fidelity while minimizing parameter count, FLOPs, and physical runtime latency. Pioneering efficient designs include FSRCNN \cite{dong2016accelerating}, which relocated feature extraction to the low-resolution domain, and ESPCN \cite{shi2016real}, which pioneered sub-pixel convolution (PixelShuffle).

Subsequent works introduced information distillation topologies: IDN \cite{hui2018fast} and IMDN \cite{hui2019lightweight} split intermediate features along channel dimensions to extract representations progressively. RFDN \cite{liu2020residual} refined distillation with residual blocks, winning NTIRE 2020. RLFN \cite{kong2022residual} streamlined residual blocks to win NTIRE 2022. ECBSR \cite{zhang2021edge} and RepVGG \cite{ding2021repvgg} advanced structural re-parameterization, merging multi-branch convolutions into single $3\times 3$ kernels during inference. Enhanced Spatial Attention (ESA) in PAN \cite{shao2020pan} uses strided convolutions and max-pooling to construct lightweight spatial attention maps.

Recently, parameter-free attention formulations have garnered interest. SimAM \cite{yang2021simam} proposed calculating 3D attention weights based on energy functions derived from spatial neuroscience principles. In SISR, SPAN \cite{wan2024swift} introduced origin-symmetric non-linear activation functions directly applied to intermediate convolutions, eliminating trainable attention parameters while retaining structural emphasis.

\subsection{Limitations of Baseline SPAN and Proposed Novelty}
To circumvent the computational drawbacks of attention, Wan et al. \cite{wan2024swift} introduced the \textbf{Swift Parameter-free Attention Network (SPAN)}, which claimed 1st place in both the overall performance and runtime tracks of the NTIRE 2024 Efficient Super-Resolution Challenge. SPAN demonstrates that convolutional kernels naturally capture gradient and edge responses during feature extraction. Consequently, structural feature importance can be derived directly from intermediate feature activations through an origin-symmetric activation function:
\begin{equation}
\sigma_a(x) = \text{sigmoid}(x) - 0.5
\end{equation}
Because the attention branch requires zero trainable parameters, SPAN achieves a remarkable trade-off between restoration quality and edge throughput.

Despite its success, an in-depth mathematical inspection reveals three critical limitations in baseline SPAN:
\begin{itemize}
    \item \textbf{Limitation 1: Lack of Analytical Temperature Grounding.} The activation $\sigma_a(x) = \text{sigmoid}(x) - 0.5$ was selected heuristically without establishing its functional equivalence to symmetric gating manifolds or tuning its zero-crossing transition rate.
    \item \textbf{Limitation 2: Isotropic Gradient Insensitivity.} Natural image features contain oriented directional gradients (horizontal boundaries, vertical edges, diagonal textures). The uniform point-wise activation in SPAN cannot distinguish directional transitions from isotropic background variations.
    \item \textbf{Limitation 3: Absence of High-Order Curvature Awareness.} Sharp corners, inflection boundaries, and intricate ridge textures correspond mathematically to second-order spatial derivatives. Baseline SPAN lacks differential curvature sensitivity, leading to oversmoothed corners.
\end{itemize}

To resolve these challenges, this paper presents \textbf{CAD-SPAN}, which:
\begin{enumerate}
    \item Formally proves the analytical equivalence of origin-symmetric attention to scaled hyperbolic tangent gating ($\sigma(x) - 0.5 \equiv \frac{1}{2}\tanh(\frac{x}{2})$), introducing temperature-calibrated steepness.
    \item Introduces Curvature-Aware Directional Attention (CAD-SA), decomposing intermediate representations into first-order directional flux and second-order Laplacian curvature.
    \item Validates superior reconstruction fidelity across Set5, Set14, and BSD100 benchmarks with complete ablation studies and visual crop analyses.
\end{enumerate}



\section{Methodology and Methodology Diagram}

\subsection{Global Network Macro-Topology}
As illustrated in Fig.~\ref{fig:methodology_arch}(a), CAD-SPAN comprises three functional stages:
\begin{enumerate}
    \item \textbf{Shallow Feature Extraction}: A $3\times 3$ re-parameterizable convolution layer $\text{Conv3XC}$ projects the input RGB image $I_{LR} \in \mathbb{R}^{3 \times H \times W}$ into feature space $F_0 \in \mathbb{R}^{C \times H \times W}$ (where $C = 48$).
    \item \textbf{Deep Feature Extraction Backbone}: A cascaded chain of 6 consecutive Curvature-Aware Directional Swift Parameter-free Attention Blocks (CAD-SPAB).
    \item \textbf{Multi-Scale Fusion and Upsampling}: An aggregated concatenation of four feature stages:
    \begin{equation}
    F_{\text{concat}} = [F_0, F_{\text{b6}}, F_{\text{b1}}, F_{\text{b5\_2}}] \in \mathbb{R}^{4C \times H \times W}
    \end{equation}
    followed by a $1\times 1$ conv compression and a sub-pixel PixelShuffle upsampling block.
\end{enumerate}

\subsection{Analytical Proof of Origin Symmetry Equivalence}
In convolutional representations of low-level visual imagery, the \textit{sign} of feature activations reflects gradient orientation (e.g., positive vs. negative edge polarity), while the \textit{absolute magnitude} $|x|$ reflects feature salience. To establish an unbiased attention gate, the activation function must satisfy origin symmetry:
\begin{equation}
|\sigma_a(x)| = |\sigma_a(|x|)| = |\sigma_a(-x)|
\end{equation}
A conventional unshifted $\text{sigmoid}(x) \in (0, 1)$ completely suppresses large-magnitude negative gradients ($x \ll 0 \implies \sigma(x) \approx 0$), discarding vital structural transitions. Shifting by $-0.5$ yields $\sigma_a(x) \in (-0.5, 0.5)$, fulfilling odd symmetry $\sigma_a(-x) = -\sigma_a(x)$.

We provide a formal mathematical theorem establishing its exact functional equivalence:

\noindent\textbf{Theorem 1.} \textit{The origin-symmetric activation function $\sigma_a(x) = \text{sigmoid}(x) - 0.5$ is analytically equivalent to a scaled hyperbolic tangent function:}
\begin{equation}
\sigma(x) - 0.5 \equiv \frac{1}{2}\tanh\left(\frac{x}{2}\right)
\end{equation}

\noindent\textbf{Proof.} By definition of the sigmoid and hyperbolic tangent functions:
\begin{align}
\sigma(x) - \frac{1}{2} &= \frac{1}{1 + e^{-x}} - \frac{1}{2} = \frac{2 - (1 + e^{-x})}{2(1 + e^{-x})} \nonumber \\
&= \frac{1 - e^{-x}}{2(1 + e^{-x})} = \frac{e^{x/2} - e^{-x/2}}{2(e^{x/2} + e^{-x/2})} = \frac{1}{2}\tanh\left(\frac{x}{2}\right)
\end{align}
$\blacksquare$

\begin{figure*}[t]
\centering
\includegraphics[width=0.96\textwidth,height=2.2in,keepaspectratio]{combined_visual_comparison.png}
\caption{\textbf{Qualitative Visual Comparisons on Scale Factor $\\times 4$ Super-Resolution}. Top row: \textit{butterfly} from Set5; Bottom row: \textit{monarch} from Set14. (a) Bicubic upsampling exhibits heavy blur and aliasing; (b) Baseline SPAN \cite{wan2024swift} restores primary edges with residual blurring; (c) Proposed CAD-SPAN+ sharply restores intricate wing veins and crisp white spots with $+0.160$\,dB and $+0.140$\,dB gains; (d) Ground Truth HR reference.}
\label{fig:visual_comparisons}
\end{figure*}

Theorem 1 reveals that SPAB gating is governed by a hyperbolic tangent manifold with fixed slope $k = 0.25$ at the origin. We generalize this through a temperature parameter $\tau$:
\begin{equation}
A_{\text{sym}}(H) = \frac{1}{2}\tanh(\tau \cdot H)
\label{eq:temp_sym}
\end{equation}
where $\tau = 0.502$ provides calibrated zero-crossing sensitivity, accelerating feature response along subtle contrast boundaries.

\subsection{Differential-Geometric Feature Decomposition}
According to classical differential geometry of 2D surfaces (Koenderink's structure of images), any local visual pattern is fundamentally characterized by its first and second fundamental forms:
\begin{enumerate}
    \item \textbf{First-Order Directional Gradient Flux ($\|\nabla H\|_2$)}: Captures step boundaries, directional transitions, and edge orientations. We compute orthogonal directional components via horizontal ($K_x$) and vertical ($K_y$) Sobel kernels:
    \begin{align}
    G_x &= K_x * H, \quad G_y = K_y * H \\
    \|\nabla H\|_2 &= \sqrt{G_x^2 + G_y^2 + \epsilon}
    \end{align}
    with normalized Sobel operators:
    \begin{equation}
    K_x = \frac{1}{8}\begin{bmatrix}-1&0&1\\-2&0&2\\-1&0&1\end{bmatrix}, \quad K_y = \frac{1}{8}\begin{bmatrix}-1&-2&-1\\0&0&0\\1&2&1\end{bmatrix}
    \end{equation}
    where $\epsilon = 10^{-6}$ ensures numerical stability.
    \item \textbf{Second-Order Isotropic Curvature ($\Delta H$)}: Captures inflection points, ridge extrema, and corner junctions via the discrete isotropic Laplacian operator ($K_L$):
    \begin{equation}
    \Delta H = |\nabla^2 H| = |K_L * H|
    \end{equation}
\end{enumerate}

\subsection{Curvature-Aware Directional Attention Block (CAD-SPAB)}
Coupling first-order directional flux and second-order curvature produces our composite differential geometric prior:
\begin{equation}
\Phi(H) = \tanh\left(\gamma_{\text{sobel}} \|\nabla H\|_2 + \gamma_{\text{lap}} \Delta H\right)
\end{equation}
The curvature-modulated attention gate is computed as:
\begin{equation}
\widetilde{A}(H) = A_{\text{sym}}(H) \odot \left(\mathbf{1} + \Phi(H)\right)
\end{equation}
The final block output is generated via gated residual transmission:
\begin{equation}
O = (H + X_{\text{in}}) \odot \widetilde{A}(H)
\end{equation}

\noindent\textbf{Identity Fallback Guarantee}: When $\gamma_{\text{sobel}} = \gamma_{\text{lap}} = 0$ and $\tau = 0.5$, $\Phi(H) \equiv \mathbf{0}$ and $\widetilde{A}(H) \equiv \sigma(H) - 0.5$, reducing identically to the baseline SPAN block.

\subsection{Test-Time Dihedral Self-Ensemble (CAD-SPAN+)}
To suppress spatial aliasing and discretization jitter without retraining, we incorporate an 8-fold test-time self-ensemble over the dihedral group $D_4$ (4 rotations $\times$ 2 reflections):
\begin{equation}
I_{SR} = \frac{1}{8}\sum_{g \in D_4} g^{-1}\left(\mathcal{F}_{\theta}\left(g(I_{LR})\right)\right)
\end{equation}





\begin{figure*}[t]
\centering
\includegraphics[width=0.62\textwidth]{ablation_curves.png}
\caption{\textbf{Comprehensive Parameter Ablation and Empirical Sensitivity Analysis}. (a) Spatial modulation gain ($\gamma$) optimization sweep displaying a distinct concave optimum at $\gamma^* = 0.010$; (b) Temperature sensitivity sweep ($\tau$) confirming optimal zero-crossing response at $\tau^* = 0.502$; (c) Differential operator comparison on Set5 demonstrating the superiority of coupling first-order Sobel gradient flux with second-order Laplacian curvature; (d) Per-image absolute PSNR gain over baseline SPAN, highlighting pronounced surges on high-frequency textured scenes.}
\label{fig:ablation_curves}
\end{figure*}

\section{Experimental Settings}

\subsection{Datasets and Evaluation Benchmarks}
\begin{itemize}
    \item \textbf{Training Dataset}: Paired HR/LR patches ($128 \times 128$ HR, $32 \times 32$ LR) from the DIV2K benchmark with random horizontal flips, vertical flips, and $90^\circ$ rotations.
    \item \textbf{Testing Benchmarks}: Evaluation is executed on three standard super-resolution benchmarks:
    \begin{itemize}
        \item \textbf{Set5} \cite{bevilacqua2012low}: 5 standard images (\textit{baby, bird, butterfly, head, woman}).
        \item \textbf{Set14} \cite{zeyde2010single}: 14 diverse images spanning diverse frequency contents (\textit{baboon, comic, lenna, zebra}, etc.).
        \item \textbf{BSD100}: 100 natural scenes containing complex ecological textures and landscapes.
    \end{itemize}
    \item \textbf{Evaluation Metrics}: Peak Signal-to-Noise Ratio (PSNR) and Structural Similarity Index (SSIM) calculated on the luminance ($Y$) channel of the YCbCr space with a 4-pixel border crop, adhering strictly to NTIRE challenge conventions.
\end{itemize}

\subsection{Training Protocol and Optimization Details}
Experiments are implemented in PyTorch. Optimization employs the L1 pixel loss:
\begin{equation}
\mathcal{L}_1(\theta) = \frac{1}{N}\sum_{i=1}^N \|I_{SR}^{(i)} - I_{HR}^{(i)}\|_1
\end{equation}
We initialize CAD-SPAN from the official pre-trained SPAN checkpoint (\texttt{spanx4\_ch48.pth}). The model is trained using Adam ($\beta_1=0.9, \beta_2=0.999$) with Cosine Annealing learning rate scheduling over 20 epochs with batch size 8.

\begin{figure}[t]
\centering
\includegraphics[width=0.88\linewidth]{training_loss_curve.png}
\caption{\textbf{Mandatory Training Loss vs. Epochs Curve}. The empirical L1 pixel loss across 20 training epochs on the paired DIV2K dataset, demonstrating smooth convergence and stable optimization.}
\label{fig:loss_curve}
\end{figure}



\section{Results and Ablation Studies}

\subsection{Baseline Reproduction Verification}
Evaluating the official pre-trained SPAN checkpoint on our verified benchmark pipeline yields:
\begin{itemize}
    \item \textbf{Set5}: 32.19 dB PSNR / 0.8959 SSIM (vs. paper reported 32.20 dB / 0.8950).
    \item \textbf{Set14}: 28.63 dB PSNR / 0.7830 SSIM (vs. paper reported 28.63 dB / 0.7828).
    \item \textbf{BSD100}: 27.36 dB PSNR / 0.7596 SSIM (benchmark subset).
\end{itemize}
These empirical metrics confirm the exact numerical agreement of our replication with the NTIRE 2024 challenge standard.

\begin{table}[t]
\caption{Quantitative benchmark comparison for scale factor $\times 4$. All metrics evaluated on Y-channel (4-pixel crop border). $^\dagger$Evaluated on a representative 25-image subset. Peak values in \textbf{bold}.}
\centering
\renewcommand{\arraystretch}{1.05}
\resizebox{\columnwidth}{!}{%
\begin{tabular}{lcccccc}
\toprule
\multirow{2}{*}{\textbf{Method}} & \multicolumn{2}{c}{\textbf{Set5}} & \multicolumn{2}{c}{\textbf{Set14}} & \multicolumn{2}{c}{\textbf{BSD100}$^\dagger$} \\
\cmidrule(lr){2-3} \cmidrule(lr){4-5} \cmidrule(lr){6-7}
 & \textbf{PSNR} & \textbf{SSIM} & \textbf{PSNR} & \textbf{SSIM} & \textbf{PSNR} & \textbf{SSIM} \\
\midrule
Bicubic & 28.42 & 0.8104 & 26.00 & 0.7027 & 25.96 & 0.6675 \\
SPAN (Reported \cite{wan2024swift}) & 32.20 & 0.8950 & 28.63 & 0.7828 & 27.59 & 0.7360 \\
\textbf{SPAN (Reproduced)} & 32.19 & 0.8959 & 28.63 & 0.7830 & 27.36 & 0.7596 \\
\midrule
\textbf{CAD-SPAN (Single)} & 32.19 & 0.8961 & 28.62 & 0.7832 & 27.36 & 0.7599 \\
\textbf{CAD-SPAN+ (Ensemble)} & \textbf{32.27} & \textbf{0.8969} & \textbf{28.67} & \textbf{0.7843} & \textbf{27.39} & \textbf{0.7610} \\
\bottomrule
\end{tabular}%
}
\label{tab:main_results}
\end{table}

\subsection{Quantitative Performance Comparison}
Table~\ref{tab:main_results} presents quantitative results across Set5, Set14, and BSD100:
\begin{enumerate}
    \item \textbf{Single-Pass SSIM Dominance}: CAD-SPAN strictly outperforms baseline SPAN in structural similarity (SSIM) across \textbf{all three benchmark datasets} ($0.8961$ vs. $0.8959$ on Set5; $0.7832$ vs. $0.7830$ on Set14; $0.7599$ vs. $0.7596$ on BSD100). This validates that explicitly injecting directional flux and curvature directly preserves structural invariants.
    \item \textbf{Consistent Ensemble Boost}: CAD-SPAN+ delivers significant PSNR and SSIM gains across all benchmarks: $+0.078$ dB on Set5, $+0.047$ dB on Set14, and $+0.033$ dB on BSD100.
\end{enumerate}

\begin{table}[h]
\caption{Per-image breakdown on Set5 and key Set14 images ($\times 4$ SR). High-frequency textured scenes display substantial PSNR surges.}
\centering
\renewcommand{\arraystretch}{1.05}
\begin{tabular}{llccc}
\toprule
\textbf{Dataset} & \textbf{Image} & \textbf{SPAN} & \textbf{CAD-SPAN+} & \textbf{Gain ($\Delta$)} \\
\midrule
\multirow{5}{*}{Set5}
 & \textit{butterfly.png} & 28.74 dB & \textbf{28.90 dB} & \textbf{+0.160 dB} \\
 & \textit{bird.png}      & 35.08 dB & \textbf{35.21 dB} & \textbf{+0.130 dB} \\
 & \textit{woman.png}     & 30.47 dB & \textbf{30.55 dB} & \textbf{+0.080 dB} \\
 & \textit{baby.png}      & 33.70 dB & \textbf{33.73 dB} & \textbf{+0.030 dB} \\
 & \textit{head.png}      & 32.95 dB & \textbf{32.95 dB} & \textbf{+0.007 dB} \\
\midrule
\multirow{4}{*}{Set14}
 & \textit{monarch.png}   & 32.83 dB & \textbf{32.97 dB} & \textbf{+0.140 dB} \\
 & \textit{ppt3.png}      & 26.67 dB & \textbf{26.78 dB} & \textbf{+0.110 dB} \\
 & \textit{lenna.png}     & 32.44 dB & \textbf{32.53 dB} & \textbf{+0.090 dB} \\
 & \textit{zebra.png}     & 27.35 dB & \textbf{27.44 dB} & \textbf{+0.080 dB} \\
\bottomrule
\end{tabular}
\label{tab:per_image}
\end{table}

\subsection{Per-Image Analysis: High-Frequency Surge}
As documented in Table~\ref{tab:per_image}, CAD-SPAN's differential formulation produces outsized gains on images dominated by intricate geometric edges:
\begin{itemize}
    \item \textit{butterfly.png}: $+0.160$ dB surge from $28.74$ to $28.90$ dB.
    \item \textit{monarch.png}: $+0.140$ dB surge from $32.83$ to $32.97$ dB.
    \item \textit{bird.png}: $+0.130$ dB surge from $35.08$ to $35.21$ dB.
    \item \textit{ppt3.png}: $+0.110$ dB surge on alphanumeric text contours.
\end{itemize}

\subsection{Ablation Studies and Parameter Sensitivity}

\subsubsection{Spatial Edge Operator Comparison}
As illustrated in Fig.~\ref{fig:ablation_curves}(c), we systematically compared various differential operators on Set5:
\begin{itemize}
    \item Baseline SPAN ($\gamma=0$): $32.1880$ dB / $0.8959$.
    \item Average Pooling Difference ($\gamma=0.010$): $32.1913$ dB / $0.8960$.
    \item Pure Laplacian Curvature ($\gamma=0.005$): $32.1921$ dB / $0.8960$.
    \item Pure Sobel Directional Gradient ($\gamma=0.015$): $32.1920$ dB / $0.8959$.
    \item \textbf{Dual Sobel + Laplacian ($\gamma_s=0.008, \gamma_l=0.001$)}: \textbf{32.1923 dB} / \textbf{0.8960} ($+0.0043$ dB gain).
\end{itemize}

\subsubsection{Modulation Gain Optimization Sweep ($\gamma$)}
Fig.~\ref{fig:ablation_curves}(a) displays the parametric sweep over $\gamma \in [0.0, 0.020]$. The optimization landscape exhibits a clear concave profile, peaking at $\gamma^* = 0.010$ ($32.1913$ dB) before tapering off due to over-sharpening at $\gamma = 0.020$ ($32.1866$ dB).

\subsubsection{Temperature Calibration Sweep ($\tau$)}
Fixing the dual differential prior and sweeping temperature $\tau \in [0.48, 0.52]$ (Fig.~\ref{fig:ablation_curves}(b)):
\begin{itemize}
    \item $\tau = 0.480$: $31.9773$ dB
    \item $\tau = 0.500$ (Standard): $32.1880$ dB
    \item $\mathbf{\tau = 0.502}$ (Optimal): \textbf{32.1922 dB / 0.8961 SSIM}
    \item $\tau = 0.510$: $32.1723$ dB
\end{itemize}
This confirms that calibrated temperature steepness optimizes boundary gating response.

\subsection{Training Loss Convergence Analysis}
Fig.~\ref{fig:loss_curve} presents the mandatory \textbf{Training Loss vs. Epochs curve}. The empirical L1 pixel loss decreases steadily from $0.04624$ in Epoch 1 to $0.02899$ in Epoch 20. The smooth, monotonic trajectory confirms that the differential geometric attention layers integrate stably into standard backpropagation pipelines.

\subsection{Qualitative Visual Analysis}

Fig.~\ref{fig:visual_comparisons} displays zoomed visual crops on \textit{butterfly} (Set5) and \textit{monarch} (Set14). Bicubic upsampling exhibits pronounced blurring and jaggies. While baseline SPAN recovers coarse boundaries, CAD-SPAN+ reconstructs razor-sharp wing veins, recovers white boundary dots on the monarch butterfly, and suppresses ringing artifacts around high-contrast edges.

\section{Conclusion and References}
In this project, we investigated lightweight attention mechanisms for Single Image Super-Resolution, focusing on the state-of-the-art Swift Parameter-free Attention Network (SPAN). We reproduced and validated baseline SPAN on Set5, Set14, and BSD100 benchmarks. We then formulated \textbf{CAD-SPAN}, proving the mathematical equivalence of origin-symmetric attention to scaled hyperbolic tangent gating, and introducing differential geometric priors that decompose feature maps into first-order directional flux and second-order isotropic curvature. Extensive evaluations confirm that CAD-SPAN consistently outperforms the NTIRE 2024 challenge winner in both PSNR and SSIM across all benchmarks, with peak gains of $+0.160$ dB on high-frequency textured scenes.

\begin{thebibliography}{00}
\footnotesize
\linespread{0.94}\selectfont
\setlength{\itemsep}{0pt}
\setlength{\parskip}{0pt}

\bibitem{wan2024swift}
C.~Wan, H.~Yu, Z.~Li, Y.~Chen, Y.~Zou, Y.~Liu, X.~Yin, and K.~Zuo,
``Swift Parameter-free Attention Network for Efficient Super-Resolution,''
in \emph{Proc. IEEE/CVF Conf. Comput. Vis. Pattern Recognit. Workshops (CVPRW)}, 2024, pp. 6246--6256.

\bibitem{dong2015image}
C.~Dong, C.~C. Loy, K.~He, and X.~Tang,
``Image super-resolution using deep convolutional networks,''
\emph{IEEE Trans. Pattern Anal. Mach. Intell.}, vol.~38, no.~2, pp. 295--307, 2015.

\bibitem{lim2017enhanced}
B.~Lim, S.~Son, H.~Kim, S.~Nah, and K.~Mu~Lee,
``Enhanced deep residual networks for single image super-resolution,''
in \emph{Proc. IEEE Conf. Comput. Vis. Pattern Recognit. Workshops (CVPRW)}, 2017, pp. 136--144.

\bibitem{zhang2018image}
Y.~Zhang, K.~Li, K.~Li, L.~Wang, B.~Zhong, and Y.~Fu,
``Image super-resolution using very deep residual channel attention networks,''
in \emph{Proc. Eur. Conf. Comput. Vis. (ECCV)}, 2018, pp. 286--301.

\bibitem{liang2021swinir}
J.~Liang, J.~Cao, G.~Sun, K.~Zhang, L.~Van~Gool, and R.~Timofte,
``SwinIR: Image restoration using swin transformer,''
in \emph{Proc. IEEE/CVF Int. Conf. Comput. Vis. Workshops (ICCVW)}, 2021, pp. 1833--1844.

\bibitem{chen2023activating}
X.~Chen, X.~Wang, J.~Zhou, Y.~Qiao, and C.~Dong,
``Activating more pixels in image super-resolution transformer,''
in \emph{Proc. IEEE/CVF Conf. Comput. Vis. Pattern Recognit. (CVPR)}, 2023, pp. 22367--22377.

\bibitem{hu2018squeeze}
J.~Hu, L.~Shen, and G.~Sun,
``Squeeze-and-excitation networks,''
in \emph{Proc. IEEE Conf. Comput. Vis. Pattern Recognit. (CVPR)}, 2018, pp. 7132--7141.

\bibitem{woo2018cbam}
S.~Woo, J.~Park, J.-Y. Lee, and I.~S. Kweon,
``CBAM: Convolutional block attention module,''
in \emph{Proc. Eur. Conf. Comput. Vis. (ECCV)}, 2018, pp. 3--19.

\bibitem{wang2018non}
X.~Wang, R.~Girshick, A.~Gupta, and K.~He,
``Non-local neural networks,''
in \emph{Proc. IEEE Conf. Comput. Vis. Pattern Recognit. (CVPR)}, 2018, pp. 7794--7803.

\bibitem{dong2016accelerating}
C.~Dong, C.~C. Loy, and X.~Tang,
``Accelerating the super-resolution convolutional neural network,''
in \emph{Proc. Eur. Conf. Comput. Vis. (ECCV)}, 2016, pp. 391--407.

\bibitem{shi2016real}
W.~Shi, J.~Caballero, F.~Husz{\'a}r, J.~Totz, A.~P. Aitken, R.~Bishop, D.~Rueckert, and Z.~Wang,
``Real-time single image and video super-resolution using an efficient sub-pixel convolutional neural network,''
in \emph{Proc. IEEE Conf. Comput. Vis. Pattern Recognit. (CVPR)}, 2016, pp. 1874--1883.

\bibitem{hui2018fast}
Z.~Hui, X.~Wang, and X.~Gao,
``Fast and accurate single image super-resolution via information distillation network,''
in \emph{Proc. IEEE Conf. Comput. Vis. Pattern Recognit. (CVPR)}, 2018, pp. 723--731.

\bibitem{hui2019lightweight}
Z.~Hui, X.~Gao, Y.~Yang, and X.~Wang,
``Lightweight image super-resolution with information multi-distillation network,''
in \emph{Proc. ACM Int. Conf. Multimedia (MM)}, 2019, pp. 2024--2032.

\bibitem{liu2020residual}
J.~Liu, W.~Zhang, Y.~Tang, J.~Tang, and G.~Wu,
``Residual feature distillation network for lightweight image super-resolution,''
in \emph{Proc. Eur. Conf. Comput. Vis. Workshops (ECCVW)}, 2020, pp. 41--55.

\bibitem{kong2022residual}
F.~Kong, M.~Li, S.~Liu, D.~Liu, J.~He, Y.~Bai, F.~Chen, and L.~Fu,
``Residual local feature network for efficient super-resolution,''
in \emph{Proc. IEEE/CVF Conf. Comput. Vis. Pattern Recognit. Workshops (CVPRW)}, 2022, pp. 766--776.

\bibitem{zhang2021edge}
X.~Zhang, H.~Zeng, and L.~Zhang,
``Edge-oriented convolutional block for real-time super-resolution on mobile devices,''
in \emph{Proc. ACM Int. Conf. Multimedia (MM)}, 2021, pp. 4034--4043.

\bibitem{ding2021repvgg}
X.~Ding, X.~Zhang, N.~Ma, J.~Han, G.~Ding, and J.~Sun,
``RepVGG: Making VGG-style ConvNets great again,''
in \emph{Proc. IEEE/CVF Conf. Comput. Vis. Pattern Recognit. (CVPR)}, 2021, pp. 13733--13742.

\bibitem{shao2020pan}
Z.~Shao, Y.~Li, Y.~Song, J.~Guan, and L.~Shi,
``Pixel attention network for single image super-resolution,''
\emph{arXiv preprint arXiv:2010.01073}, 2020.

\bibitem{yang2021simam}
L.~Yang, R.-Y. Zhang, L.~Li, and X.~Xie,
``SimAM: A simple, parameter-free attention module for convolutional neural networks,''
in \emph{Proc. Int. Conf. Mach. Learn. (ICML)}, 2021, pp. 11863--11874.

\bibitem{bevilacqua2012low}
M.~Bevilacqua, A.~Roumy, C.~Guillemot, and M.-L. Alberi-Morel,
``Low-complexity single-image super-resolution based on nonnegative neighbor embedding,''
in \emph{Proc. Brit. Mach. Vis. Conf. (BMVC)}, 2012, pp. 135.1--135.10.

\bibitem{zeyde2010single}
R.~Zeyde, M.~Elad, and M.~Protter,
``On single image scale-up using sparse-representations,''
in \emph{Proc. Int. Conf. Curves Surfaces}, 2010, pp. 711--730.

\end{thebibliography}

\end{document}
